\begin{table}[H]
\centering
\caption{Nama Daerah Perbandingan Cluster Biplot vs Statis (2024)} \label{tab:appbiplot_statis_name_2024}
\begin{adjustbox}{max width=\textwidth}
\begin{tabular}{|p{4cm}|p{10cm}|}
\hline
\textbf{Kategori} & \textbf{Daftar Wilayah} \\
\hline
Statis 1 (Rendah) \& Biplot 1 (Rendah) & Cianjur, Tasikmalaya, Ciamis, Majalengka, Pangandaran, Pandeglang, Lebak \\
\hline
Statis 1 (Rendah) \& Biplot 2 (Sedang) & - \\
\hline
Statis 1 (Rendah) \& Biplot 3 (Tinggi) & - \\
\hline
Statis 2 (Sedang) \& Biplot 1 (Rendah) & Kep. Seribu, Garut, Indramayu, Kota Banjar, Serang \\
\hline
Statis 2 (Sedang) \& Biplot 2 (Sedang) & Bogor, Sukabumi, Bandung, Kuningan, Cirebon, Sumedang, Subang, Purwakarta, Karawang, Bandung Barat, Kota Cirebon, Kota Tasikmalaya, Tangerang, Kota Cilegon, Kota Serang \\
\hline
Statis 2 (Sedang) \& Biplot 3 (Tinggi) & - \\
\hline
Statis 3 (Tinggi) \& Biplot 1 (Rendah) & - \\
\hline
Statis 3 (Tinggi) \& Biplot 2 (Sedang) & - \\
\hline
Statis 3 (Tinggi) \& Biplot 3 (Tinggi) & Kota Jakarta Selatan, Kota Jakarta Timur, Kota Jakarta Pusat, Kota Jakarta Barat, Kota Jakarta Utara, Bekasi, Kota Bogor, Kota Sukabumi, Kota Bandung, Kota Bekasi, Kota Depok, Kota Cimahi, Kota Tangerang, Kota Tangerang Selatan \\
\hline
\end{tabular}
\end{adjustbox}
\end{table}